% !TEX root = ../tfm.tex
\chapter{Adaptación de MEG-XL a ``EEG-XL''}
\providecommand{\codepath}[1]{\nolinkurl{#1}}

\section{Descripción general de la propuesta}

Como se ha ido comentando a lo largo de los anteriores capítulos, la propuesta de este trabajo consiste en adaptar MEG-XL a registros EEG. La idea es mantener, dentro de lo que sea posible, aquello que hace valioso al modelo original, como por ejemplo el aprendizaje autosupervisado con contexto largo. MEG-XL no se limita a procesar ventanas breves alineadas con palabras, sino que aprende representaciones a partir de fragmentos continuos de dos minutos y medio. Esta propiedad resulta especialmente interesante para EEG, donde la señal es más ruidosa y la información lingüística puede aparecer distribuida en patrones débiles a lo largo del tiempo.

La Figura~\ref{fig:megxl_overview} resume el flujo original de MEG-XL que se toma como punto de partida: primero un preentrenamiento autosupervisado sobre señal MEG continua y después un ajuste fino supervisado para clasificación contextual de palabras.

\begin{figure}[htbp]
    \raggedright
    \resizebox{0.99\textwidth}{!}{%
    \begin{tikzpicture}[
        font=\footnotesize,
        node distance=5mm,
        box/.style={draw, rounded corners=1.5mm, minimum width=2.45cm,
                    minimum height=1.15cm, align=center, thick, fill=blue!7},
        model/.style={draw, rounded corners=1.5mm, minimum width=2.75cm,
                    minimum height=1.3cm, align=center, thick, fill=orange!10},
        loss/.style={draw, rounded corners=1.5mm, minimum width=2.6cm,
                    minimum height=1.15cm, align=center, thick, fill=green!10},
        arrow/.style={-{Latex[length=2.4mm]}, thick}
    ]
        \node[font=\bfseries\large] (pretitle) {Preentrenamiento autosupervisado};
        \node[box, below=5mm of pretitle] (meg) {MEG continuo\\2,5 min a 50 Hz};
        \node[box, right=of meg] (bio) {BioCodec\\tokens por canal};
        \node[box, right=of bio] (mask) {Enmascaramiento\\de bloques temporales};
        \node[model, right=of mask] (cc) {Transformer\\criss--cross};
        \node[loss, right=of cc] (mtp) {Predicción de\\tokens ocultos};
        \draw[arrow] (meg) -- (bio);
        \draw[arrow] (bio) -- (mask);
        \draw[arrow] (mask) -- (cc);
        \draw[arrow] (cc) -- (mtp);

        \node[font=\bfseries\large, below=22mm of pretitle] (fintitle) {Ajuste fino para decodificación de palabras};
        \node[box, below=5mm of fintitle] (wordmeg) {Ventanas alineadas\\con palabras};
        \node[model, right=of wordmeg] (pretrained) {MEG-XL\\preentrenado};
        \node[model, right=of pretrained] (seq) {Contextualización\\de la secuencia};
        \node[loss, right=of seq] (t5) {D-SigLIP\\contra embeddings T5};
        \node[loss, right=of t5] (topk) {Recuperación\\Top-$k$};
        \draw[arrow] (wordmeg) -- (pretrained);
        \draw[arrow] (pretrained) -- (seq);
        \draw[arrow] (seq) -- (t5);
        \draw[arrow] (t5) -- (topk);
        \draw[arrow, dashed, bend right=18] (cc.south) to (pretrained.north);

        \node[draw, dashed, rounded corners, fit=(meg)(mtp), inner sep=3mm] {};
        \node[draw, dashed, rounded corners, fit=(wordmeg)(topk), inner sep=3mm] {};
    \end{tikzpicture}}
    \caption{Resumen del preentrenamiento y del ajuste fino de MEG-XL. El modelo aprende primero a reconstruir tokens cerebrales en contextos largos y posteriormente reutiliza esas representaciones en una tarea contrastiva de clasificación contextual de palabras. Adaptado conceptualmente de Jayalath y Parker Jones.}
    \label{fig:megxl_overview}
\end{figure}


Esta adaptación se puede entender como un cambio de dominio. El modelo de partida fue entrenado con MEG, utilizando sensores magnéticos, geometrías específicas y una relación señal-ruido distinta. En EEG cambian la magnitud física, la referencia, la disposición de los sensores, el número de canales y los artefactos más frecuentes. Por tanto, es necesario construir una interfaz común que permita representar diferentes montajes EEG dentro de la misma arquitectura.

El entrenamiento final lanzado en este trabajo corresponde a un preentrenamiento continuo en dos etapas. Primero se utilizan datasets de lectura y después datasets de escucha. La etapa de lectura sirve como una adaptación inicial a EEG lingüístico en condiciones más controladas. La etapa de escucha aproxima el modelo al dominio final, que es la clasificación de palabras durante percepción auditiva continua.

Además, el experimento compara tres variantes preentrenadas y un control supervisado directo. Una variante aprende el conjunto EEG desde una inicialización aleatoria. Las otras dos parten del checkpoint de MEG-XL y se diferencian únicamente en la forma de representar el tipo de sensor: una utiliza un embedding específico de EEG inicializado a partir del embedding de magnetómetro, y la otra reutiliza directamente el embedding de magnetómetro para los canales EEG. Las tres variantes preentrenadas siguen exactamente el mismo orden, lectura seguida de escucha, de modo que la comparación permite estudiar por separado el efecto del preentrenamiento EEG, de la transferencia desde MEG y de la elección del embedding de tipo de sensor.

La propuesta final se puede resumir en una arquitectura que recoge cuatro decisiones de MEG-XL:

\begin{itemize}
\setlength{\itemsep}{0pt}
\setlength{\parsep}{0pt}
\setlength{\topsep}{2pt}
\setlength{\partopsep}{0pt}
\item utilizar ventanas continuas de 150 segundos;
\item filtrar inicialmente la señal entre 0,1 y 40 Hz en la frecuencia de muestreo original y remuestrearla después a 50 Hz, de modo que la banda representada en la entrada final queda limitada por la frecuencia de Nyquist final;
\item tokenizar cada canal con BioCodec y entrenar mediante predicción de tokens enmascarados;
\item realizar el entrenamiento en dos fases, primero lectura y después escucha, comparando inicialización aleatoria, transferencia desde MEG-XL con embedding EEG y transferencia desde MEG-XL con embedding MEG.
\end{itemize}

De esta forma, la parte supervisada de clasificación de palabras queda separada del preentrenamiento, como sucede con MEG-XL. Los cuatro datasets descritos en el capítulo anterior se utilizan para intentar aprender una representación ``general'' de EEG. Posteriormente, ds004408 permite ajustar esa representación a una tarea \textit{word-aligned} concreta, con ventanas asociadas al inicio de cada palabra y evaluación mediante recuperación de embeddings lingüísticos.

La Figura~\ref{fig:eegxl_design} muestra cómo se materializa esta comparación experimental en EEG-XL, separando el preentrenamiento EEG desde cero, las dos variantes de transferencia desde MEG-XL y el control supervisado directo.

\begin{figure}[htbp]
    \raggedright
    \resizebox{0.98\textwidth}{!}{%
    \begin{tikzpicture}[
        font=\small,
        arrow/.style={-{Latex[length=3mm]}, thick},
        init/.style={draw, rounded corners=2mm, minimum width=4.2cm,
                    minimum height=1.22cm, align=center, thick, fill=orange!10},
        curriculum/.style={draw, rounded corners=2mm, minimum width=5.4cm,
                          minimum height=1.95cm, align=center, thick, fill=blue!7},
        final/.style={draw, rounded corners=2mm, minimum width=4.9cm,
                     minimum height=1.35cm, align=center, thick, fill=green!10},
        control/.style={draw, rounded corners=2mm, minimum width=4.9cm,
                       minimum height=1.35cm, align=center, thick, fill=gray!10},
        note/.style={font=\footnotesize, align=center}
    ]
        \node[init] (scratchinit) at (0,5.4) {Inicialización aleatoria\\arquitectura EEG-XL};
        \node[curriculum] (scratchcurr) at (6.2,5.4) {\textbf{Currículo EEG autosupervisado}\\
              lectura: Nieuwland + ZuCo 2.0\\
              $\downarrow$ mejor checkpoint\\
              escucha: SparrKULee + ds007808};
        \node[final] (scratchft) at (12.7,5.4) {Fine-tuning en ds004408\\conjunto EEG desde cero};

        \node[init] (megeeginit) at (0,2.2) {Inicialización desde MEG-XL\\pesos compatibles};
        \node[curriculum] (megeegcurr) at (6.2,2.2) {\textbf{Currículo EEG autosupervisado}\\
              embedding EEG específico\\
              lectura $\rightarrow$ escucha\\
              mismo conjunto y particiones};
        \node[final] (megeegft) at (12.7,2.2) {Fine-tuning en ds004408\\transferencia MEG-XL\\con embedding EEG};

        \node[init] (megmeginit) at (0,-1.0) {Inicialización desde MEG-XL\\pesos compatibles};
        \node[curriculum] (megmegcurr) at (6.2,-1.0) {\textbf{Currículo EEG autosupervisado}\\
              embedding MEG reutilizado\\
              lectura $\rightarrow$ escucha\\
              mismo conjunto y particiones};
        \node[final] (megmegft) at (12.7,-1.0) {Fine-tuning en ds004408\\transferencia MEG-XL\\con embedding MEG};

        \node[control] (directinit) at (0,-4.0) {Inicialización aleatoria};
        \node[control] (direct) at (6.2,-4.0) {Sin currículo autosupervisado};
        \node[control] (directft) at (12.7,-4.0) {Fine-tuning directo en ds004408\\control supervisado};

        \draw[arrow] (scratchinit) -- (scratchcurr);
        \draw[arrow] (scratchcurr) -- (scratchft);
        \draw[arrow] (megeeginit) -- (megeegcurr);
        \draw[arrow] (megeegcurr) -- (megeegft);
        \draw[arrow] (megmeginit) -- (megmegcurr);
        \draw[arrow] (megmegcurr) -- (megmegft);
        \draw[arrow] (directinit) -- (direct);
        \draw[arrow] (direct) -- (directft);

        \node[note] at (6.2,-5.45) {Las cuatro condiciones comparten el mismo conjunto de fine-tuning y el mismo protocolo de evaluación.};
    \end{tikzpicture}}
    \caption{Diseño experimental de la adaptación de MEG-XL a EEG-XL. Las tres rutas autosupervisadas mantienen el mismo currículo de lectura--escucha y difieren en la inicialización y en el embedding de tipo de sensor; el control supervisado permite separar el efecto del preentrenamiento EEG de la evaluación directa sobre ds004408.}
    \label{fig:eegxl_design}
\end{figure}


\section{Preprocesamiento del EEG}

\subsection{Preprocesado de la señal}

Todos los datasets de entrenamiento se transforman a una frecuencia de muestreo común de 50 Hz tal y como lo hace MEG-XL. Esta elección reduce el coste computacional de las ventanas largas. Un segmento de 150 segundos contiene 7500 muestras por canal antes de la tokenización, por lo que aumentar la frecuencia de muestreo incrementaría de forma directa la memoria necesaria y el número de tokens que debe procesar el Transformer.

El entrenamiento realizado aplica primero un filtrado paso banda de 0,1--40 Hz sobre la señal en su frecuencia de muestreo original y, posteriormente, remuestrea todos los registros a 50 Hz. Como la frecuencia de Nyquist de una señal muestreada a 50 Hz es 25 Hz, la entrada final del modelo no debe interpretarse como una representación completa de la banda 0,1--40 Hz. En la práctica, durante el remuestreo se aplica el filtrado antialiasing necesario, por lo que las componentes superiores a la frecuencia de Nyquist final quedan atenuadas o eliminadas. Por tanto, aunque el preprocesamiento inicial utiliza un límite superior de 40 Hz, la banda efectiva disponible para el modelo tras el remuestreo queda limitada aproximadamente a 0,1--25 Hz.

Se mantiene la frecuencia final de 50 Hz por compatibilidad con la configuración original de MEG-XL y por coste computacional. En ventanas de 150 segundos, esta frecuencia produce 7500 muestras por canal antes de la tokenización. Aumentar la frecuencia de muestreo incrementaría directamente la longitud temporal, el número de tokens y el consumo de memoria del Transformer, dificultando el entrenamiento con contextos largos.

Después del filtrado y el remuestreo, cada grabación se divide en ventanas completas de 150 segundos. Para montar estas ventanas se utilizan segmentos continuos y no solapados dentro de cada registro. Los fragmentos finales que no alcanzan la duración requerida se descartan. Esta decisión evita introducir relleno artificial, repetir muestras para completar una ventana o unir partes de la señal que no son contiguas, lo que podría introducir artefactos y sesgos en el entrenamiento.

A partir de cada ventana de 150 segundos, se subdivide la señal en bloques de tres segundos. Esta subdivisión mantiene la escala temporal utilizada por MEG-XL y coincide con la duración de las ventanas que más tarde se utilizan en clasificación de palabras. Durante el preentrenamiento continuo, esos bloques no se asocian a palabras concretas; funcionan como unidades temporales sobre las que se puede aplicar el enmascaramiento.

La normalización se aplica para reducir diferencias de escala entre canales, sesiones y datasets. Este paso es de vital importancia en EEG porque cada conjunto puede haber sido registrado con amplificadores, referencias y montajes distintos. La normalización no elimina por completo estas diferencias, pero evita que el modelo dependa de magnitudes absolutas que no son comparables entre datasets.

Siguiendo el ejemplo de MEG-XL, se ha creado una interfaz común para cargar los datasets. Cada ejemplo contiene la señal EEG, las posiciones de los sensores y una máscara que indica qué canales son reales. Cuando un conjunto tiene menos sensores que otro, no se fuerza una interpolación espacial previa; se representa dentro de una dimensión común y se enmascaran los canales ausentes. Así, el modelo puede entrenar con diferentes montajes sin confundir sensores inexistentes con señales de valor cero.

\begin{figure}[htbp]
    \raggedright
    \resizebox{0.96\textwidth}{!}{%
    \begin{tikzpicture}[
        font=\small,
        node distance=8mm,
        arrow/.style={-{Latex[length=3mm]}, thick},
        box/.style={draw, rounded corners=2mm, minimum width=3.7cm,
                   minimum height=1.55cm, align=center, thick, fill=blue!7},
        adapt/.style={draw, rounded corners=2mm, minimum width=3.8cm,
                     minimum height=1.65cm, align=center, thick, fill=orange!10},
        output/.style={draw, rounded corners=2mm, minimum width=5.0cm,
                      minimum height=2.0cm, align=center, thick, fill=green!10},
        smallnote/.style={font=\footnotesize, align=center}
    ]
        \node[box] (raw) {EEG heterogéneo\\frecuencias\\referencias\\canales};
        \node[box, right=of raw] (filter) {Paso banda\\$0{,}1$--$40\,\mathrm{Hz}$};
        \node[box, right=of filter] (resample) {Remuestreo\\$50\,\mathrm{Hz}$\\antialias};
        \node[box, right=of resample] (norm) {Normalización\\escala};

        \node[adapt, below=12mm of norm] (windows) {Ventanas\\$150\,\mathrm{s}$\\$7500$ muestras};
        \node[adapt, left=of windows] (blocks) {50 bloques\\$3\,\mathrm{s}$\\máscara temporal};
        \node[adapt, left=of blocks] (montage) {Montaje común\\$C_{\max}$\\relleno};

        \node[output, below=12mm of blocks] (sample) {Entrada\\
        $\mathbf{X}\in\mathbb{R}^{C_{\max}\times 7500}$,\quad
        $\mathbf{P}\in\mathbb{R}^{C_{\max}\times 3}$\\
        $\mathbf{m}_{\mathrm{sensor}}\in\{0,1\}^{C_{\max}}$,\quad
        $s_c=2$\\
        máscara temporal};

        \draw[arrow] (raw) -- (filter);
        \draw[arrow] (filter) -- (resample);
        \draw[arrow] (resample) -- (norm);
        \draw[arrow] (norm) -- (windows);
        \draw[arrow] (windows) -- (blocks);
        \draw[arrow] (blocks) -- (montage);
        \draw[arrow] (montage.south) |- (sample.west);

        \node[smallnote, below=3mm of windows] {restos $<150\,\mathrm{s}$: descarte};
    \end{tikzpicture}}
    \caption{Preprocesamiento EEG: filtrado, remuestreo, normalización, ventanas, montaje común y máscaras.}
    \label{fig:eegxl_preprocessing}
\end{figure}


\subsection{Alineamiento con las palabras}

El entrenamiento autosupervisado final no utiliza palabras, transcripciones ni alineaciones lingüísticas como etiquetas. Durante las etapas de lectura y escucha, el modelo recibe EEG continuo y aprende a reconstruir códigos neuronales ocultos. Por tanto, pueden combinarse datasets con idiomas y paradigmas diferentes sin construir un vocabulario común.

En los conjuntos de lectura, las anotaciones de presentación visual o fijación sirven para seleccionar tareas relacionadas con procesamiento lingüístico, pero no se utilizan como objetivo de entrenamiento. En ZuCo se emplea la tarea de lectura natural, y en el conjunto de Nieuwland et al. las tareas \textit{delong} y \textit{control}. En los conjuntos de escucha, se seleccionan tareas de percepción auditiva continua. En OpenNeuro ds007808, los registros \textit{listeningcovert} se tratan con cuidado: la grabación completa puede proporcionar contexto, pero solo los intervalos etiquetados como escucha son candidatos para ser enmascarados y reconstruidos.

El alineamiento palabra a palabra aparece después, en la etapa de \textit{fine-tuning} con OpenNeuro ds004408. En ese caso sí se extraen ventanas de tres segundos alrededor del inicio de cada palabra y se agrupan 50 palabras consecutivas para formar una entrada de 150 segundos. De esta forma primero el modelo aprende regularidades generales de EEG continuo y después ya se especializa en recuperación de palabras.

\section{Adaptación de la arquitectura MEG-XL}

\subsection{Componentes conservados}

La adaptación final mantiene la arquitectura principal de MEG-XL. La señal se tokeniza con BioCodec, que transforma cada canal en una secuencia discreta de códigos mediante cuantización vectorial residual. En la configuración utilizada, BioCodec emplea un vocabulario de 256 códigos, seis cuantizadores y un factor de reducción temporal de 12. De esta forma, los 7500 puntos temporales de una ventana de 150 segundos a 50 Hz se convierten en una secuencia de tokens más corta, manejable por el Transformer.

Sobre esos tokens se aplica el Transformer con una combinación de 2 tipos de atención que se llama atención cruzada \textit{criss-cross}. La atención temporal permite relacionar momentos dentro de un mismo sensor, mientras que la atención espacial combina información entre sensores para un mismo instante. Esta factorización conserva la idea original de MEG-XL: modelar dependencias largas sin construir una atención completa sobre todos los pares canal-tiempo.

También se conserva el objetivo autosupervisado. El modelo oculta bloques temporales y debe predecir los códigos de BioCodec correspondientes a esas posiciones. La \textit{loss} es una entropía cruzada sobre los códigos discretos de los cuantizadores. En el entrenamiento se enmascaran subsegmentos de tres segundos, de forma que el modelo no puede resolver la tarea únicamente interpolando muestras vecinas inmediatas y fuerza al Transformer a aprender patrones más amplios de la señal.

Por tanto, la contribución de la adaptación está en hacer que la misma arquitectura pueda recibir EEG heterogéneo y en definir un conjunto de entrenamiento adecuado para pasar de lectura a escucha. La arquitectura, el tokenizer y el objetivo autosupervisado se mantienen lo más cerca posible de MEG-XL, mientras que los cambios se concentran en la entrada EEG, las máscaras de sensores, la selección de datasets y el entrenamiento por etapas.

\subsection{Cambios para trabajar con canales EEG}

MEG-XL fue diseñado para trabajar con magnetómetros y gradiómetros. En EEG solo existe un tipo de sensor dentro del experimento, pero los montajes varían mucho entre datasets. Por ello se introduce el tipo de sensor EEG dentro del espacio de tipos ya utilizado por la arquitectura. La configuración final mantiene tres tipos posibles: gradiómetro, magnetómetro y EEG. En los datos de este trabajo todos los canales reales pertenecen al tipo EEG, pero conservar los otros tipos permite cargar parcialmente los pesos de MEG-XL sin rediseñar el modelo.

La posición de cada electrodo se utiliza como información espacial al no tener el índice de un canal el mismo significado en todos los conjuntos. Un canal situado en una región frontal no debería tratarse como equivalente a otro canal occipital solo porque ocupa la misma posición en una matriz. Al proporcionar coordenadas y máscaras de sensores, el modelo puede aprender relaciones espaciales en función de la geometría real disponible en cada registro.

Formalmente, si el canal \(c\) tiene coordenadas normalizadas \(\mathbf{p}_c=(x_c,y_c,z_c)\in\mathbb{R}^3\) y tipo de sensor \(s_c\in\{0,1,2\}\), con \(0\) para gradiómetro, \(1\) para magnetómetro y \(2\) para EEG, la información espacial se introduce mediante un \textit{embedding} Fourier gaussiano:

\begin{equation}
\boldsymbol{\phi}(\mathbf{p}_c)=
\left[
\sin\left(2\pi \mathbf{p}_c \mathbf{B}\right),
\cos\left(2\pi \mathbf{p}_c \mathbf{B}\right)
\right],
\qquad
\mathbf{B}_{ij}\sim\mathcal{N}(0,\sigma^2),
\end{equation}

\noindent donde \(\mathbf{B}\in\mathbb{R}^{3\times d_F/2}\) es una matriz de frecuencias fija. Este vector se proyecta a la dimensión latente del modelo:

\begin{equation}
\mathbf{e}^{\text{pos}}_c =
\mathbf{W}_{\text{pos}}\boldsymbol{\phi}(\mathbf{p}_c)
+ \mathbf{b}_{\text{pos}}.
\end{equation}

El tipo de sensor se representa mediante una tabla aprendida \(\mathbf{E}_{\text{type}}\in\mathbb{R}^{3\times d}\):

\begin{equation}
\mathbf{e}^{\text{type}}_c =
\mathbf{E}_{\text{type}}[s_c].
\end{equation}

Por tanto, para el token temporal \(\tau\) del canal \(c\), el vector de entrada al Transformer queda:

\begin{equation}
\mathbf{h}^{(0)}_{c,\tau} =
\mathbf{e}^{\text{tok}}_{c,\tau}
+ \mathbf{e}^{\text{pos}}_c
+ \mathbf{e}^{\text{type}}_c.
\end{equation}

En la configuración base de EEG se usa \(s_c=2\) para todos los canales reales. La posición distingue electrodos dentro de un mismo montaje y el tipo de sensor permite mantener la compatibilidad con los pesos de MEG-XL sin confundir la modalidad física de la señal. En la variante de transferencia con embedding MEG, la señal sigue tratándose como EEG en posiciones, máscaras y orientación, pero la consulta a la tabla de tipos reutiliza la fila de magnetómetro aprendida por MEG-XL. Así se puede aislar el efecto de introducir un embedding EEG específico frente a reutilizar directamente el embedding MEG disponible.

Otra diferencia con respecto a MEG-XL es el tamaño de nuestro modelo debido al número de canales máximo de los datasets. Nuestros datos de EEG contienen montajes de 64, 105, 128 o configuraciones variables según el laboratorio y la sesión. La dimensión máxima de canales se infiere a partir de los datasets disponibles y a cada uno se le aplica una máscara. Esta máscara se aplica tanto durante el procesamiento como en el cálculo de la pérdida, evitando que los canales añadidos para igualar dimensiones modifiquen canales entrenados.

\begin{table}[htbp]
\centering
\caption{Resumen de la arquitectura autosupervisada EEG-XL y número de parámetros del codificador.}
\label{tab:arquitectura_modelo_eeg}
\footnotesize
\setlength{\tabcolsep}{5pt}
\renewcommand{\arraystretch}{1.15}
\begin{tabularx}{\textwidth}{
    c
    >{\raggedright\arraybackslash}p{4.2cm}
    >{\raggedright\arraybackslash}X
    >{\centering\arraybackslash}p{1.4cm}
}
\toprule
\textbf{N.} &
\textbf{Componente} &
\textbf{Tipo} &
\textbf{Parámetros} \\
\midrule

0 &
\texttt{tokenizer} &
\texttt{BioCodecTokenizerAdapter} &
3,2 M \\

1 &
\texttt{rvq\_projector} &
\texttt{Linear} &
49,7 K \\

2 &
\texttt{position\_fourier\_emb} &
\texttt{GaussianFourierEmb3D} &
375 \\

3 &
\texttt{position\_projector} &
\texttt{Linear} &
128 K \\

4 &
\texttt{orientation\_fourier\_emb} &
\texttt{GaussianFourierEmb3D} &
375 \\

5 &
\texttt{orientation\_projector} &
\texttt{Linear} &
128 K \\

6 &
\texttt{sensor\_type\_layer} &
\texttt{Embedding} &
1,5 K \\

7 &
\texttt{criss\_cross\_transformer} &
\texttt{SpatialTemporalEncoder} &
21,0 M \\

8 &
\texttt{output\_head} &
\texttt{Linear} &
787 K \\

\midrule
\multicolumn{3}{r}{\textbf{Parámetros entrenables}} &
22,1 M \\

\multicolumn{3}{r}{\textbf{Parámetros no entrenables}} &
3,2 M \\

\multicolumn{3}{r}{\textbf{Parámetros totales}} &
25,3 M \\

\bottomrule
\end{tabularx}
\end{table}

Dos experimentos utilizan como base el modelo preentrenado con MEG de MEG-XL. En ambos casos, el modelo se carga en la nueva arquitectura de forma parcial. Los pesos compatibles, como los del Transformer y gran parte de las proyecciones internas, pueden reutilizarse. Los componentes que dependen de la nueva configuración EEG o que dejan de encajar exactamente se inicializan de nuevo. La diferencia entre ambas variantes está en la fila de embedding de tipo de sensor empleada para los canales EEG: una introduce un embedding EEG específico y la otra reutiliza el embedding de magnetómetro. De esta forma, la carga parcial permite estudiar la transferencia de MEG a EEG sin forzar una equivalencia falsa entre sensores de diferentes tareas.

\begin{figure}[htbp]
    \raggedright
    \resizebox{0.97\textwidth}{!}{%
    \begin{tikzpicture}[
        font=\small,
        arrow/.style={-{Latex[length=3mm]}, thick},
        kept/.style={draw, rounded corners=2mm, minimum width=3.55cm,
                    minimum height=1.45cm, align=center, thick, fill=blue!7},
        changed/.style={draw, rounded corners=2mm, minimum width=3.55cm,
                       minimum height=1.45cm, align=center, thick, fill=orange!10},
        objective/.style={draw, rounded corners=2mm, minimum width=3.65cm,
                         minimum height=1.45cm, align=center, thick, fill=green!10},
        note/.style={font=\footnotesize, align=center}
    ]
        \node[kept] (input) {EEG\\$[B,C,7500]$};
        \node[kept, right=8mm of input] (codec) {BioCodec\\congelado\\$Q=6$, $V=256$, $r=12$};
        \node[kept, right=8mm of codec] (codes) {Códigos RVQ\\$[B,C,625,6]$};
        \node[kept, right=8mm of codes] (proj) {Proyección RVQ\\$\mathbf{e}^{\mathrm{tok}}_{c,\tau}\in\mathbb{R}^{512}$};

        \node[changed, below=11mm of codec] (pos) {Posición 3D\\Fourier\\$\mathbf{e}^{\mathrm{pos}}_c$};
        \node[changed, right=8mm of pos] (type) {Tipo\\EEG: $s_c=2$\\$\mathbf{e}^{\mathrm{type}}_c$};
        \node[changed, left=8mm of pos] (ori) {Orientación\\MEG activa\\EEG cero};

        \node[changed, below=11mm of type, minimum width=6.1cm, minimum height=1.75cm] (sum) {Embedding de entrada\\
        $\mathbf{h}^{(0)}_{c,\tau}=\mathbf{e}^{\mathrm{tok}}_{c,\tau}+\mathbf{e}^{\mathrm{pos}}_c+\mathbf{e}^{\mathrm{type}}_c$};
        \node[changed, left=10mm of sum] (mask) {Enmascaramiento\\bloques $3\,\mathrm{s}$\\válidos};

        \node[kept, right=13mm of sum, minimum width=5.8cm, minimum height=2.6cm] (cc) {\textbf{Transformer criss-cross $\times 8$}\\[1mm]
        $d/2$: espacial\\
        $d/2$: temporal\\
        concat. + FFN};
        \node[objective, right=11mm of cc] (head) {Salida\\$512\rightarrow Q\times V$\\logits};
        \node[objective, below=8mm of head] (loss) {Entropía cruzada\\tokens ocultos\\$Q=6$};

        \draw[arrow] (input) -- (codec);
        \draw[arrow] (codec) -- (codes);
        \draw[arrow] (codes) -- (proj);
        \draw[arrow] (proj.south) to[out=-90,in=35] (sum.north east);
        \draw[arrow] (pos) -- (sum);
        \draw[arrow] (type) -- (sum);
        \draw[arrow] (ori) -- (sum);
        \draw[arrow] (mask) -- (sum);
        \draw[arrow] (sum) -- (cc);
        \draw[arrow] (cc) -- (head);
        \draw[arrow] (head) -- (loss);

        \node[note, below=6mm of mask, text width=5.0cm] {Azul: MEG-XL.\quad Naranja: EEG.};
    \end{tikzpicture}}
    \caption{Arquitectura EEG-XL: BioCodec, códigos RVQ, embeddings, Transformer criss-cross y predicción de tokens enmascarados.}
    \label{fig:eegxl_architecture}
\end{figure}


\subsection{Tokenización de la señal}

Durante el desarrollo se contemplaron otros tokenizadores, como BrainOmni o BrainTokenizer, porque están diseñados para representar bioseñales heterogéneas y podrían ser adecuados para combinar modalidades. Sin embargo, finalmente el entrenamiento utiliza BioCodec. Así se mantiene la compatibilidad con MEG-XL y se reduce el número de cambios simultáneos: si se cambiara a la vez la modalidad, el tokenizer y la arquitectura, sería más difícil interpretar el efecto de la transferencia. Por otro lado, es un tokenizador entrenado sobre EEG, por lo que, teóricamente, debería funcionar mejor con nuestro conjunto de datos que con el de MEG-XL.

BioCodec se aplica canal a canal. Cada canal EEG se convierte en una secuencia de códigos discretos y el Transformer aprende sobre esos códigos junto con la información de posición y tipo de sensor. El uso de seis cuantizadores permite representar la señal mediante varios niveles discretos. La salida del modelo predice, para cada posición enmascarada, el código correspondiente en cada cuantizador.

\subsection{Modelado temporal y espacial}

La arquitectura final trabaja con segmentos largos de 150 segundos. Esta escala es mucho mayor que la utilizada por la mayoría de modelos EEG centrados en ventanas de pocos segundos. El objetivo es que las capas del Transformer aprendan tanto patrones locales como dependencias más amplias relacionadas con el estado de la grabación, la continuidad del estímulo y la dinámica temporal del procesamiento lingüístico.

La atención \textit{criss-cross} separa explícitamente dos dimensiones. En la dimensión temporal, el modelo analiza cómo evoluciona la señal dentro de un sensor. En la dimensión espacial, integra información entre sensores manteniendo el mismo instante tokenizado. Esta separación reduce el coste respecto a una atención completa y se adapta de forma natural a la representación canal-tiempo del EEG.

En los registros \textit{listeningcovert} de ds007808 se añade una restricción adicional mediante una máscara de objetivo. El contexto puede incluir toda la grabación, pero la pérdida solo se aplica sobre intervalos de escucha. Esto evita que el modelo sea entrenado para reconstruir como objetivo principal periodos de habla encubierta, que no corresponden directamente a la percepción auditiva que se quiere transferir a ds004408.

\section{Entrenamiento del modelo}

\subsection{Preentrenamiento autosupervisado}

Cada entrenamiento ejecuta dos fases consecutivas:

\begin{enumerate}
\setlength{\itemsep}{0pt}
\setlength{\parsep}{0pt}
\setlength{\topsep}{2pt}
\setlength{\partopsep}{0pt}
\item entrenamiento continuo con datasets de lectura;
\item entrenamiento continuo con datasets de escucha, inicializado desde el mejor checkpoint de la fase de lectura correspondiente.
\end{enumerate}

La fase de lectura utiliza Nieuwland et al. y ZuCo 2.0 en lectura natural. La fase de escucha utiliza SparrKULee y OpenNeuro ds007808 (JapanEEG). La validación respeta separaciones por sujeto o por sesión según el dataset, como se describió en el capítulo anterior.

El orden de lectura $\rightarrow$ escucha se escogió para que los datasets de lectura proporcionen una primera adaptación a EEG lingüístico con estímulos visuales y condiciones relativamente controladas. Después, los datasets de escucha ajustan la representación al dominio auditivo continuo, que es el más próximo a la evaluación final con ds004408. El checkpoint producido por lectura se pasa explícitamente a la etapa de escucha, por lo que no se reutilizan checkpoints de experimentos anteriores de escucha aislada.

De esta forma el entrenamiento del modelo tiene la siguiente configuración: BioCodec, frecuencia final de 50 Hz, banda 0,1--40 Hz, ventanas de 150 segundos, batch size 4 en lectura y batch size 1 en escucha, validación cada 500 pasos y guardado periódico de checkpoints. El entrenamiento se ejecuta con una GPU por variante preentrenada, de forma que las configuraciones pueden lanzarse en paralelo pero no comparten pesos entre sí.

En MEG-XL se realizó el entrenamiento con un batch size de 1. En la fase de lectura, el menor número de canales de EEG permite aumentar el lote a 4 sin superar la memoria de una GPU A100 de 40 GB. En la fase de escucha se mantiene batch size 1, porque los registros auditivos y la configuración de canales hacen que el coste de memoria sea mayor. Por tanto, la comparación entre configuraciones no se basa en cambiar el tamaño de lote entre inicializaciones: todas las variantes preentrenadas utilizan batch size 4 en lectura y batch size 1 en escucha.

\subsection{Clasificación contextual de palabras}

El modelo de MEG-XL destaca por tener pocos parámetros entrenables para poder adaptarse a tareas supervisadas concretas. En este trabajo se mantiene la misma filosofía: el preentrenamiento autosupervisado aprende representaciones generales de EEG y la evaluación final se realiza mediante una tarea supervisada de clasificación contextual de palabras.

La clasificación contextual de palabras no forma parte del preentrenamiento autosupervisado continuo, sino que constituye la tarea supervisada utilizada para evaluar la calidad de las representaciones aprendidas. Una vez completadas las etapas de lectura y escucha, los checkpoints resultantes se emplean para inicializar el \textit{fine-tuning} sobre OpenNeuro ds004408. Este conjunto se reserva para esta fase y no se incorpora al conjunto de datos previo, evitando que el modelo haya observado previamente durante el preentrenamiento los registros utilizados en la evaluación final.

Las anotaciones temporales de ds004408 se obtienen de los niveles de palabras incluidos en los ficheros \textit{TextGrid}, descartando los niveles correspondientes a fonemas. Para cada palabra se extrae una ventana EEG de tres segundos, comprendida entre 0,5 segundos antes de su inicio y 2,5 segundos después. Posteriormente, se concatenan las ventanas correspondientes a 50 palabras consecutivas, produciendo una entrada de 150 segundos. Por tanto, la dimensión temporal coincide con la empleada durante el preentrenamiento, aunque en este caso no se trata necesariamente de 150 segundos continuos de la grabación original, sino de una secuencia estructurada de ventanas alineadas con palabras.

El modelo procesa conjuntamente las 50 ventanas y genera una representación para cada palabra. Sobre las características producidas por el Transformer se aplica una cabeza de proyección que transforma cada representación cerebral en un vector de 1024 dimensiones. Como objetivo lingüístico se utilizan embeddings extraídos de la capa intermedia 12 de \textit{T5-large}. El aprendizaje se realiza mediante una pérdida contrastiva de tipo SigLIP, que aproxima la representación predicha al embedding de la palabra correcta y la separa de los embeddings correspondientes al resto de palabras presentes en el lote.

Durante la evaluación, cada embedding cerebral predicho se compara mediante similitud coseno con los embeddings lingüísticos de un conjunto de recuperación. Se consideran vocabularios formados por las 50 y las 250 palabras más frecuentes y se calcula la precisión \textit{Top-10}, junto con su variante balanceada por palabra. Esta evaluación sigue la formulación de recuperación semántica utilizada en los trabajos de clasificación contextual de palabras, aunque los resultados solo pueden compararse directamente cuando coinciden el conjunto de candidatos, la partición de los datos y el protocolo de evaluación.

% El experimento final compara tres condiciones. La primera entrena la arquitectura de clasificación desde una inicialización aleatoria sobre ds004408. La segunda parte del checkpoint obtenido mediante el conjunto autosupervisado de lectura y escucha cuya arquitectura EEG fue inicializada desde cero. La tercera utiliza el checkpoint del mismo conjunto, pero iniciado a partir de MEG-XL. La comparación entre la primera y la segunda configuración permite medir la contribución del preentrenamiento autosupervisado con EEG, mientras que la comparación entre la segunda y la tercera permite aislar la aportación adicional de la transferencia desde MEG. De esta manera, puede determinarse si las mejoras proceden del conjunto EEG, de los pesos originales de MEG-XL o de la combinación de ambos factores.

\begin{figure}[htbp]
    \centering
    \resizebox{0.96\textwidth}{!}{%
    \begin{tikzpicture}[
        font=\small,
        arrow/.style={-{Latex[length=3mm]}, thick},
        box/.style={draw, rounded corners=2mm, minimum width=3.05cm,
                   minimum height=1.4cm, align=center, thick, fill=blue!7},
        language/.style={draw, rounded corners=2mm, minimum width=3.7cm,
                        minimum height=1.45cm, align=center, thick, fill=orange!10},
        train/.style={draw, rounded corners=2mm, minimum width=4.15cm,
                     minimum height=1.55cm, align=center, thick, fill=green!10},
        note/.style={font=\footnotesize, align=center}
    ]
        \node[box] (textgrid) at (0,3.6) {TextGrid\\$t_i$};
        \node[box] (epoch) at (3.6,3.6) {Ventana EEG\\$t_i-0{,}5:t_i+2{,}5$\\$3\,\mathrm{s}$};
        \node[box] (sequence) at (7.2,3.6) {50 palabras\\$150\,\mathrm{s}$};
        \node[box] (encoder) at (10.8,3.6) {EEG-XL\\contexto};
        \node[box] (project) at (14.4,3.6) {Proyección\\$\widehat{\mathbf{y}}_i\in\mathbb{R}^{1024}$};

        \draw[arrow] (textgrid) -- (epoch);
        \draw[arrow] (epoch) -- (sequence);
        \draw[arrow] (sequence) -- (encoder);
        \draw[arrow] (encoder) -- (project);

        \node[language] (t5) at (5.4,0.8) {T5-large\\capa 12\\$\mathbf{y}_i\in\mathbb{R}^{1024}$};
        \node[train] (siglip) at (10.8,0.8) {\textbf{Ajuste}\\SigLIP\\pares correctos/incorrectos};

        \node[language] (vocab) at (5.4,-1.9) {Vocabulario\\50/250 palabras\\embeddings T5};
        \node[train] (retrieval) at (10.8,-1.9) {\textbf{Evaluación}\\coseno + ranking\\Top-10};

        \draw[arrow] (t5) -- (siglip);
        \coordinate (split) at ($(project.south)+(0,-8mm)$);
        \draw[thick] (project.south) -- (split);
        \draw[arrow] (split) -| (siglip.north);
        \draw[arrow] (vocab) -- (retrieval);
        \draw[arrow, dashed] (split) -- ++(8mm,0) |- (retrieval.east);
    \end{tikzpicture}}
    \caption{Ajuste y evaluación de palabras: ventanas alineadas, contexto EEG-XL, proyección cerebral, objetivo T5 y recuperación por ranking.}
    \label{fig:eegxl_word_finetuning}
\end{figure}

